\section{Discussion}
\label{sec:discussion}

\subsection{Water level cannot separate these configurations}

The six members differ in wave coupling, bed mobility and the spatial
treatment of vegetation resistance, and they produce water level records that
are almost the same. Anomaly RMSE varies by six millimetres or less between
members at any station, and no member is consistently best across the three
stations. Over the same simulations Liu--Weisberg skill spans 0.505 to 0.655
and endpoint separation spans 334 to 661~m.

A calibration exercise conducted on water level alone would therefore have no
basis on which to prefer any of these configurations, and would plausibly
select whichever member happened to minimise a station-averaged RMSE by a
margin smaller than the measurement uncertainty. This is the case for treating
Lagrangian observations as a distinct constraint rather than as a confirmation
of a model already tuned to tide gauges. It also suggests that the frequent
practice of reporting coastal model skill through water level statistics alone
may certify transport behaviour that was never examined.

The reason is structural rather than incidental. Water level in a micro-tidal
basin with two open connections is set largely by the offshore forcing and the
basin's capacity to fill and drain, both of which the six members share.
Transport depends on the velocity field, which they do not.

\subsection{Wave coupling}

Wave coupling reduces Lagrangian skill by 0.035 to 0.042, a small effect but a
consistent one, negative in ten of the twelve deployments. It was included in
this model because spectral waves generated in the western Sicilian channel
reach both inlets, and because canopy attenuation is a documented process in
Mediterranean seagrass systems. The result does not contradict either premise.
It says that in this basin, over this window, representing waves did not
improve the prediction of surface transport, and slightly degraded it.

The penalty does not depend on how vegetation resistance is represented. Measured
separately on the two roughness axes it is 0.035 and 0.042, statistically
indistinguishable, with the same consistency across deployments and the same two
deployments improving in both. A single-axis measurement could not have
established that.

One limit remains. The window contains one moderate swell event rather than a
storm, so the test is of ordinary summer conditions, not of the conditions under
which wave effects on transport would be expected to be largest. The wave
contribution is also measured on a fixed bed only, because the mobile-bed
no-wave configurations could not be run, for reasons taken up below.

\subsection{What the interaction acts on}

Neither bed mobility nor distributed vegetation roughness changes Lagrangian
skill by more than a marginal amount on its own: 0.019 for bed mobility under
uniform roughness, 0.001 for roughness on a fixed bed with waves, 0.009 without
them. Applied together they give 0.131 and 0.149 depending on which is held
fixed, improving 32 of the 35 drifters in the first case, and produce the only
configuration that reproduces the observed transport rate, at 0.97 of the
observed path length against 0.78 to 0.93 for the others. The joint effect
exceeds the sum of the separate point estimates by an order of magnitude, which
is the signature of an interaction.

Three candidate explanations for it were tested and rejected before a fourth
line of evidence narrowed the question.

The first was that the roughness field organises where the bed erodes,
producing a channelised bathymetry that conveys flow more efficiently. The
bed-change fields do not support it. Their magnitudes are nearly identical
between the two mobile-bed members, median $-0.059$ against $-0.047$~m in the
lagoon, and their spatial difference is close to zero across the interior. The share of
total bed change held by the busiest 5\% of lagoon cells is 18.9\% and 19.0\%,
indistinguishable, so neither is more focused than the other.

The second was that canopy drag suppresses erosion beneath the meadows. Both
mobile-bed members show more bed change under the canopy than on bare
substrate, in a ratio of 1.49 under uniform roughness and 1.57 under
distributed roughness, so the effect runs in the wrong direction and the
roughness treatment makes it slightly stronger rather than weaker.

The third was that the successful member simply carries a stronger current
field. It does not. Its lagoon-mean surface speed is 0.078~\si{\metre\per\second},
below that of a member whose particles travel 20\% too slowly. The vertical profiles show that
bed mobility, not roughness, controls the shear, and that adding canopy drag to
a mobile bed lowers surface speed rather than raising it.

The particles nevertheless travel at the right speed. Since it is not the mean
speed of the field, it has to be where and when the particles encounter the
faster water, which is a property of the joint space and time structure of the
flow rather than of its magnitude.

Stating this openly is preferable to the available alternative, which is to
attribute the improvement to canopy physics because the successful member is
the one with a canopy. None of the three diagnostics supports that.

The canopy parametrisation does behave as intended, which makes the negative
result sharper rather than weaker. Holding the bed fixed and moving from
uniform to distributed roughness attenuates the wave field in proportion to the
vegetation present, reducing mean orbital velocity at the bed by 14.3\% over
\emph{Posidonia}, 13.7\% over \emph{Cymodocea}, 11.3\% over sand and 3.3\% over
the exposed rock substrate near the inlets. The ordering follows the canopy
density, and the same treatment slows the near-bed flow over \emph{Posidonia}
roughly three times as much as over sand. The vegetation is therefore doing to
the waves and to the bed layer what a canopy should. What it does not do, on a
fixed bed, is change surface transport, and that is the gap the interaction
occupies.

Splitting the trajectory error into a magnitude part and a direction part
narrows where the interaction acts (Figure~\ref{fig:decomp}). For each drifter
the simulated path length, the simulated net displacement and the mean angular
difference between simulated and observed step directions were computed over
the scored window.

The four fixed-bed members are indistinguishable in all three. Path length runs
0.77 to 0.80 of observed, net displacement 0.78 to 0.82, and mean heading error
15.8 to 16.2$^\circ$. Distributed roughness on a fixed bed leaves the heading
error at 16.0$^\circ$, unchanged to the first decimal, so it does not alter
where the particles go any more than it alters how far.

Bed mobility acts mainly on magnitude, raising path length to 0.94 and net
displacement to 0.95, with a smaller gain in heading, to 12.5$^\circ$. Adding
distributed roughness to a mobile bed acts mainly on direction, cutting the
heading error to 8.4$^\circ$, a third of what remained, for a further 0.04 in
path length. The roughness treatment therefore has a directional effect that
exists only when the bed can move, and that is the interaction stated in terms
of what it corrects. Figure~\ref{fig:tracks} shows the same distinction on
individual tracks.

This is a narrowing, not a full account. Treating each track as a single vector
and reconstructing the endpoint separation from the mean displacement ratio and
the mean heading error reproduces the four fixed-bed members within 12\%, at
507 to 538~m against 549 to 583~m observed, but underpredicts the two
mobile-bed members by a factor of 1.4 to 1.9, at 354 and 239~m against 661 and
334~m. A single mean heading is therefore an incomplete description of the
mobile-bed error, and what remains unexplained is why.

\begin{figure}[H]
  \centering
  \includegraphics[width=\textwidth]{transport_error_decomposition}
  \caption{The trajectory error split into magnitude and direction. (a) Path
  length and (b) net displacement, both as the ratio of simulated to observed.
  (c) Mean angular difference between simulated and observed step directions.
  The four fixed-bed members coincide in all three panels.}
  \label{fig:decomp}
\end{figure}

\subsection{A mobile bed does not stabilise without waves}

The two ensemble cells combining a mobile bed with no wave coupling could not be
obtained. Both abort on the velocity cap, reaching 27 and
43~\si{\metre\per\second} against a limit of 25, in a single offshore cell and
after roughly half a day of simulation. The four wave-coupled members, carrying
the same sediment settings and the same cap in the earlier configuration, never
approached it.

The behaviour survived the obvious explanation. Morphodynamics initially ran
over the entire domain, including water hundreds of metres deep, and the first
aborts occurred in a cell 27~m deep, so confining sediment to depths above 20~m
was the natural fix. It did not work. With the restriction in place both members
aborted again, at the same cell, at comparable velocities. Restricting the bed
did allow the two wave-coupled mobile-bed members to complete, so the
restriction is not itself the problem.

What this says is that wave coupling performs stabilising work in a
morphodynamic simulation of this basin, and that removing it exposes a feedback
the wave field was otherwise damping. A plausible route is through bed shear
stress, to which waves contribute directly and which sets the erosion rate, so
that removing the wave contribution alters the balance between erosion and the
flow field that drives it. This work does not establish that mechanism, and the
result is reported as an observed property of the configuration.

It also bounds the ensemble. Wave coupling and bed mobility cannot be varied
independently here, so what the design supports is a statement about wave
coupling on a fixed bed and about the roughness-by-bed-mobility interaction
under wave coupling, not a full three-way attribution.

\subsection{The bed-mobility caveat}

The configuration that reproduces observed transport depends on morphodynamics
whose critical shear stress for erosion has not been constrained by a repeated
bathymetric survey. Over nine days the lagoon bed changes by up to 1.0~m in
isolated cells, against a median of 0.05~m, in a basin whose mean depth is
0.95~m. A metre of bed change in nine days is not a physically plausible rate
for this system. The morphodynamic component therefore functions in this
ensemble as a perturbation to the bed whose realism is unverified, and the
interaction reported above should be read with that in mind.

Two readings remain open. Bed mobility may be standing in for a deficiency in
the initial bathymetry, in which case the improvement is a compensation and
would not survive a corrected bed. Or the model may genuinely require an
evolving bed to reproduce transport in a system where the seagrass canopy and
the sediment surface adjust to one another, in which case the calibration
problem is real but the physics is too. Distinguishing them requires
observational constraint on the bed, and satellite-derived bathymetry over
repeated epochs is the most tractable route currently available.

\subsection{Vertical structure}

Surface speed exceeds bed speed by a factor between 1.9 and 4.5 in the lagoon
interior, in a basin whose mean depth is below one metre. Even the weakest case
places a factor of nearly two between the layer a surface drifter samples and
the layer that governs bed shear stress. A depth-averaged representation of
this basin would assign a single velocity to both.

This is consistent with what \citet{DeMarchis2012} found from field
measurements and three-dimensional modelling in the same lagoon, and it
supplies a quantitative basis for the choice of a three-dimensional solver
here. The dependence of the shear ratio on bed mobility, from near 2.0 with a
fixed bed to near 4.0 with a mobile one, also indicates that the vertical
structure is not an independent property of the basin but responds to the same
model choices that govern transport.

